\section{Conclusion}
\label{sec:conclusion}
AI-assisted visual analytics derives its value from the computational contribution to a scientific task and the judgments scientists can inspect and revise. Similar interfaces can distribute work differently: configuring a pipeline differs from teaching a model, and manipulating a representation differs from specifying a generative constraint.

The three perspectives work together as a map, a lens, and a light. The Assistance $\times$ Visualization Modality map locates approaches, the task lens makes their analytical purposes comparable, and the light of analytic agency reveals how consequential choices are distributed between scientists and computation. Read together, these perspectives shift attention from the presence of AI or an immersive interface to the work a mechanism enables, the judgments it influences, and the opportunities it leaves for inspection and correction. The central question is how a particular division of work supports a scientific purpose, and what evidence establishes that support. Progress from hairballs to hypotheses is meaningful when scientists can explain not only what an assisted analysis suggests, but why its evidence warrants the next question.
